\section{Introduction}
\label{sec:intro}

Large language models (LLMs) are increasingly embedded in decision support, education, and automation pipelines. Even when a deployment is not ``safety-critical'' in the narrow sense, unreliable outputs, unsafe instructions, privacy failures, or biased generalizations can cause material harm. National guidance therefore treats AI risk management as an organizational discipline, not a single metric \citep{nist2023airmf}.

This document evaluates \textbf{one} target LLM configuration using a fixed questionnaire of fourteen items supplied for Software Verification and Validation (Spring 2026, Quiz \#3). Each item is tied to one of seven \textbf{trustworthy AI characteristics}: (1)~Valid and Reliable, (2)~Safe, (3)~Secure and Resilient, (4)~Explainable and Interpretable, (5)~Privacy-Enhanced, (6)~Fair with Harmful Bias Managed, and (7)~Accountable and Transparent.

\paragraph{Research stance.}
Rather than treating the quiz as disconnected short answers, we adopt an explicit empirical template:
\begin{enumerate}
  \item \textbf{Research question} --- what compliance behavior is under test?
  \item \textbf{Hypothesis} --- predicted rubric outcome and mechanism (e.g., refusal, partial explanation, or failure mode).
  \item \textbf{Experiment} --- exact prompt, model settings, and logging procedure.
  \item \textbf{Findings} --- Compliant (C), Partially Compliant (P), or Non-Compliant (N), with a short justification anchored in the model's text.
\end{enumerate}

\paragraph{Deliverables mapped to the course brief.}
The repository supplies (i)~Python source to replay the prompt battery and archive raw completions as JSON, and (ii)~this PDF, which records rubric scores and narrative analysis. A separate Canvas screenshot with your name visible is still required by the grader; we document how to produce it in Appendix~\ref{app:canvas}.

\paragraph{Reading guide.}
Section~\ref{sec:prelim} summarizes the NIST framing and rubric. Section~\ref{sec:method} states the measurement protocol. Section~\ref{sec:experiments} contains the hypotheses, experiments, and results for all fourteen items, followed by a consolidated score table.
